\documentclass{article}
\usepackage[utf8]{inputenc}
\usepackage{amsmath,amssymb}
\usepackage[most]{tcolorbox}
\usepackage{geometry}
\geometry{margin=2.5cm}

\newcommand{\step}[1]{\par\noindent\hangindent=3.5em\hangafter=1%
  \makebox[3.5em][l]{\textbf{#1.}}}
\newcommand{\steptag}[1]{\hfill{\small\textsf{[#1]}}}

\newtcolorbox{proofbox}[1]{
  colback=gray!5, colframe=gray!60,
  fonttitle=\bfseries, title={#1},
  breakable, enhanced,
  left=4pt, right=4pt, top=4pt, bottom=4pt,
  before skip=10pt, after skip=10pt
}

\begin{document}

\begin{proofbox}{Parameterized bound quotient H-space package: for every u...}
\step{1} Parameterized bound quotient H-space package: for every universal def-stage1-polar-witness-data tuple W=(C\_rect,C\_ch,C\_pol,C\_grp,C\_path,C\_der,e\_rect,kappa\_ch,kappa\_pol,kappa\_der,delta\_*,epsilon\_*\textasciicircum{}r,e\_S1,r\_iso) with C\_rect,C\_ch,C\_pol,C\_grp,C\_path,C\_der>=1, 0<e\_rect<=1/C\_rect, 0<kappa\_ch,kappa\_pol,kappa\_der<=1/2, delta\_*=min\{1/4,kappa\_ch/(4*C\_ch),kappa\_pol/(4*C\_pol)\}, epsilon\_*\textasciicircum{}r=min\{1/4,kappa\_ch/(4*C\_ch),kappa\_pol/(4*C\_pol),kappa\_der/(8*C\_der),1/C\_grp,delta\_*/(12*C\_path*C\_grp)\}, e\_S1=min\{e\_rect,epsilon\_*\textasciicircum{}r/C\_rect\}, and r\_iso=min\{delta\_*/4,kappa\_der/(8*C\_der)\}, there exist universal e\_quot\textasciicircum{}r>0 and epsilon\_B\textasciicircum{}r>0 with epsilon\_B\textasciicircum{}r=min\{epsilon\_*\textasciicircum{}r,e\_quot\textasciicircum{}r\} such that, for every finite-dimensional exact-unit epsilon\_r-C*-algebra with 0<=epsilon\_r<=epsilon\_B\textasciicircum{}r and 1<dim\_C calX<infinity, and for displayed data satisfying all of the following: (A\_2) for every V in calUbar\_delta\_* and A\textasciicircum{}par in B\_\{2delta\_*\}\textasciicircum{}\{icalH\}(0) there is a unique g\_V(A\textasciicircum{}par) in B\_\{2delta\_*\}\textasciicircum{}\{calH\}(0) with f\_V(A\textasciicircum{}par+g\_V(A\textasciicircum{}par))=0, where f\_V(A)=(1/2)*(((J+A\textasciicircum{}dagger) bold-dot V\textasciicircum{}dagger) bold-dot (V bold-dot (J+A))-J), chi\_V(A\textasciicircum{}par)=V bold-dot (J+A\textasciicircum{}par+g\_V(A\textasciicircum{}par)) lies in calU, ||Dg\_V(A\textasciicircum{}par)||<=C\_ch*(epsilon\_r+delta\_*), ||D\_\{A\textasciicircum{}perp\}f\_V(A\textasciicircum{}par+g\_V(A\textasciicircum{}par))-I\_calH||<=C\_ch*(epsilon\_r+delta\_*)<1, and these C\textasciicircum{}1 charts cover calU; (A\_4) set Pi\_delta\_*:calU x B\_\{delta\_*\}\textasciicircum{}\{calH\}(J)->calX by Pi\_delta\_*(U,K):=U bold-dot K and set S\_delta\_*:=Pi\_delta\_*(calU x B\_\{delta\_*\}\textasciicircum{}\{calH\}(J)); Pi\_delta\_* is a C\textasciicircum{}1 diffeomorphism onto S\_delta\_* with unique inverse maps u\_delta:S\_delta\_*->calU and h\_delta:S\_delta\_*->B\_\{delta\_*\}\textasciicircum{}\{calH\}(J) satisfying X=u\_delta(X) bold-dot h\_delta(X) for every X in S\_delta\_* and u\_delta(U)=U and h\_delta(U)=J for every U in calU; (A\_5) the displayed maps mu:calU x calU->calU and sigma:calU->calU are defined by mu(U,V):=u\_delta(U bold-dot V) and sigma(U):=u\_delta(U\textasciicircum{}dagger), are global C\textasciicircum{}1 and, for every U,V,Z in calU, mu(J,U)=mu(U,J)=U, sigma(J)=J, ||mu(U,V)-U bold-dot V||<=C\_grp*epsilon\_r, ||sigma(U)-U\textasciicircum{}dagger||<=C\_grp*epsilon\_r, ||mu(mu(U,V),Z)-mu(U,mu(V,Z))||<=C\_grp*epsilon\_r, ||mu(sigma(U),U)-J||<=C\_grp*epsilon\_r, and ||mu(U,sigma(U))-J||<=C\_grp*epsilon\_r; (A\_6) for every U\_0,U\_1 in calU and q in [0,1] satisfying ||U\_1-U\_0||<=q, C\_path*q<=1/4, and C\_path*(q+epsilon\_r*q+q\textasciicircum{}2)<delta\_*-C\_pol*(epsilon\_r*delta\_*+delta\_*\textasciicircum{}2), the path H(-,U\_0,U\_1):[0,1]->calU given by H(t,U\_0,U\_1):=u\_delta((1-t)*U\_0+t*U\_1) is defined, is jointly continuous in its displayed variables, and joins U\_0 to U\_1, and satisfies H(t,cU\_0,cU\_1)=c*H(t,U\_0,U\_1) for every c in U(1) and t in [0,1]; and (R) set q\_*:=C\_grp*epsilon\_r, set r\_-:=delta\_*-C\_pol*(epsilon\_r*delta\_*+delta\_*\textasciicircum{}2), and set eta\_*:=C\_path*(q\_*+epsilon\_r*q\_*+q\_*\textasciicircum{}2); the guards C\_ch*(epsilon\_r+delta\_*)<=kappa\_ch, C\_pol*(epsilon\_r+delta\_*)<=kappa\_pol, q\_*<r\_-, C\_path*q\_*<=1/4, and eta\_*<r\_- hold; then, for those same u\_delta,h\_delta,mu,sigma,H, there exist a space breve-calU and maps breve-mu:breve-calU x breve-calU->breve-calU and breve-sigma:breve-calU->breve-calU such that breve-calU=calU\_e/U(1), set breve-e:=[J], breve-mu([U],[V])=[mu(U,V)], breve-sigma([U])=[sigma(U)], (breve-calU,breve-mu,breve-e) is a connected H-space, breve-sigma is a smooth left inversion, sigma(cU)=conj(c)*sigma(U), and ||sigma(U)-U\textasciicircum{}dagger||<=C\_grp*epsilon\_r for every c in U(1) and U in calU, and breve-calU is a connected compact orientable smooth manifold without boundary of real dimension (dim\_C calX)-1. \steptag{validated} \\[6pt]
\step{1.1} Fix an arbitrary witness tuple W and arbitrary displayed algebra/data satisfying the hypotheses. By positivity of every entry in the defining minimum, epsilon\_*\textasciicircum{}r>0. Invoke lem-stage1-quotient-manifold-package and let e\_quot\textasciicircum{}r>0 be its universal constant; set epsilon\_B\textasciicircum{}r:=min\{epsilon\_*\textasciicircum{}r,e\_quot\textasciicircum{}r\}>0. Since epsilon\_r<=epsilon\_B\textasciicircum{}r<=e\_quot\textasciicircum{}r and 1<dim\_C calX<infinity, that external gives breve-calU:=calU\_e/U(1) as a connected compact orientable smooth manifold without boundary of real dimension dim\_C(calX)-1. \steptag{validated} \\[4pt]
\step{1.2} Under (A\_2), (A\_4), (A\_5), and (R), the registered smooth-upgrade externals give a smooth embedded structure on calU, the same smooth polar inverse (u\_delta,h\_delta), and smooth maps alpha(c,U)=cU, mu(U,V)=u\_delta(U bold-dot V), sigma(U)=u\_delta(U\textasciicircum{}dagger), with unchanged point values and derivatives, satisfying mu(cU,dV)=cd mu(U,V), sigma(cU)=conj(c)sigma(U), and the bound ||sigma(U)-U\textasciicircum{}dagger||<=C\_grp*epsilon\_r. \steptag{validated} \\[6pt]
\step{1.2.1} For each graph point in (A\_2), (R) gives ||D\_\{A\textasciicircum{}perp\}f\_V-I\_calH||<=C\_ch*(epsilon\_r+delta\_*)<=kappa\_ch<=1/2<1. The Neumann-series criterion therefore makes D\_\{A\textasciicircum{}perp\}f\_V invertible. All remaining hypotheses of lem-stage1-smooth-unitary-atlas are exactly (A\_2), so that external upgrades the same g\_V and chi\_V, without changing their points or first derivatives, to a C\textasciicircum{}infinity graph atlas defining a smooth embedded manifold structure on calU. \steptag{validated} \\[4pt]
\step{1.2.2} With the smooth atlas from the preceding step, (A\_4) says that Pi\_delta\_* is a bijective C\textasciicircum{}1 diffeomorphism, hence a bijective C\textasciicircum{}1 local diffeomorphism; its image S\_delta\_* is open because every local diffeomorphism is locally open. Thus lem-stage1-smooth-polar-inverse applies and upgrades the same ambient-bilinear Pi\_delta\_* and its same set-theoretic inverse (u\_delta,h\_delta), without changing points or first derivatives, to smooth maps. \steptag{validated} \\[6pt]
\step{1.2.2.1} Write N:=dim\_C(calX).  The conjugate-linear involution gives the real direct-sum decomposition calX\_R=calH direct-sum icalH: for X in calX, X=(X+X\textasciicircum{}dagger)/2+(X-X\textasciicircum{}dagger)/2 with the two summands respectively Hermitian and anti-Hermitian, their intersection is zero, and multiplication by i is a real-linear isomorphism calH->icalH.  Hence dim\_R(calH)=dim\_R(icalH)=N and dim\_R(calX)=2N.  The smooth graph charts supplied by node 1.2.1 are modeled on open subsets of icalH, so dim\_R(calU)=N; consequently calU x B\_\{delta\_*\}\textasciicircum{}\{calH\}(J) and the ambient real manifold calX both have dimension 2N. \steptag{validated} \\[4pt]
\step{1.2.2.2} By (A\_4), Pi\_delta\_* is a C\textasciicircum{}1 diffeomorphism of the 2N-dimensional domain onto its image S\_delta\_*; in particular it is a C\textasciicircum{}1 immersion, so at every (U,K) its ambient differential D Pi\_delta\_*(U,K):T\_U calU direct-sum calH -> calX\_R is injective.  By the equal 2N-dimensional source and target count in the preceding child, this ambient differential is a linear isomorphism.  The finite-dimensional inverse-function theorem therefore supplies, at every (U,K), a domain neighborhood O\_\{U,K\} whose image Pi\_delta\_*(O\_\{U,K\}) is open in ambient calX and on which Pi\_delta\_* is a C\textasciicircum{}1 diffeomorphism.  These images cover S\_delta\_*, so S\_delta\_* is ambient-open and Pi\_delta\_* is a bijective C\textasciicircum{}1 local diffeomorphism as a map into calX, not merely relative to S\_delta\_*. \steptag{validated} \\[4pt]
\step{1.2.2.3} Now all hypotheses of lem-stage1-smooth-polar-inverse hold: node 1.2.1 supplies the smooth embedded atlas, (A\_4) supplies the same ambient-bilinear bijection Pi\_delta\_* and the same set-theoretic inverse (u\_delta,h\_delta), and the preceding child supplies the required ambient openness and local-diffeomorphism property.  That external therefore upgrades Pi\_delta\_* and (u\_delta,h\_delta) to smooth maps without changing their point values or C\textasciicircum{}1 differentials, which is the conclusion asserted by node 1.2.2. \steptag{validated} \\[4pt]
\step{1.2.3} In (A\_5), global well-typedness of u\_delta(U bold-dot V) and u\_delta(U\textasciicircum{}dagger), whose domain is S\_delta\_*, entails U bold-dot V,U\textasciicircum{}dagger in S\_delta\_* for all U,V in calU. Apply lem-stage1-explicit-smooth-unitary-operations to the graph family and atlas of the first step and the polar maps of the second: the resulting maps have precisely alpha(c,U)=cU, mu(U,V)=u\_delta(U bold-dot V), sigma(U)=u\_delta(U\textasciicircum{}dagger), are smooth with unchanged C\textasciicircum{}1 differentials, and obey mu(cU,dV)=cd mu(U,V) and sigma(cU)=conj(c)sigma(U). The pointwise estimate ||sigma(U)-U\textasciicircum{}dagger||<=C\_grp*epsilon\_r is the unchanged (A\_5) estimate. \steptag{validated} \\[4pt]
\step{1.3} On the component calU\_e, the operations mu and sigma descend to smooth maps breve-mu([U],[V]):=[mu(U,V)] and breve-sigma([U]):=[sigma(U)]; with breve-e=[J], the exact unit identities and the (A\_6) path applied at q\_*=C\_grp*epsilon\_r make (breve-calU,breve-mu,breve-e) a connected H-space and breve-sigma a smooth left inversion. Together with the manifold conclusion and the retained equivariance/error estimate, these are all conclusions of node 1. \steptag{validated} \\[6pt]
\step{1.3.1} Because calU\_e is a connected component of the smooth manifold calU, it is open and hence a smooth submanifold. The connected set U(1) x calU\_e has connected image under alpha and contains alpha(1,J)=J, so alpha preserves calU\_e. The restricted action is free because every U in calU has a right inverse and is nonzero, so cU=U forces c=1. It is proper: for Phi(c,U)=(cU,U) and compact K in calU\_e x calU\_e, Phi\textasciicircum{}\{-1\}(K) is closed and lies in the compact set U(1) x pr\_2(K). Lem-topology-quotient-manifold therefore equips breve-calU:=calU\_e/U(1), with its quotient topology, with the unique smooth structure making p:calU\_e->breve-calU a smooth surjective submersion; choose this structure, without identifying it with the package atlas. The package manifold is the same orbit space with the same quotient topology. Thus connectedness and compactness persist. If n=dim\_C calX-1, the package orientation is a coherent choice of generators of H\_n(breve-calU,breve-calU minus \{x\};Z). Local homology and its continuation maps depend only on the topology. In a smooth m-manifold, the local half-space calculation says that nonzero local homology at every point occurs precisely in degree m and precisely at interior points. Hence the chosen quotient structure has m=n, no boundary, and inherits the same coherent generators as an orientation. Continuity of mu and sigma shows that their images of calU\_e x calU\_e and calU\_e are connected, contain J, and stay in calU\_e. Equivariance gives p(mu(cU,dV))=p(mu(U,V)) and p(sigma(cU))=p(sigma(U)), so breve-mu([U],[V]):=[mu(U,V)] and breve-sigma([U]):=[sigma(U)] are well-defined. Local smooth sections of p x p and p express them as p o mu o s\_12 and p o sigma o s, proving smoothness. Finally (A\_5) gives breve-mu(breve-e,[U])=[U]=breve-mu([U],breve-e) and breve-sigma(breve-e)=breve-e. \steptag{validated} \\[6pt]
\step{1.3.1.1} Because calU\_e is a connected component of the smooth manifold calU, it is open and hence a smooth submanifold. The connected set U(1) x calU\_e has connected image under alpha and that image contains alpha(1,J)=J, so alpha preserves calU\_e. This restricted action is free: U in calU has a right inverse, hence U is nonzero, and cU=U forces c=1. It is proper: for Phi(c,U)=(cU,U), if K is compact in calU\_e x calU\_e, then Phi\textasciicircum{}\{-1\}(K) is closed, since K is closed in the Hausdorff manifold product, and is contained in the compact set U(1) x pr\_2(K); hence Phi\textasciicircum{}\{-1\}(K) is compact. Lem-topology-quotient-manifold therefore equips the orbit space breve-calU:=calU\_e/U(1), with its quotient topology, with a unique smooth structure for which p:calU\_e->breve-calU is a smooth surjective submersion. We choose this smooth structure. \steptag{archived} \\[4pt]
\step{1.3.1.2} The smooth manifold asserted by lem-stage1-quotient-manifold-package is the same orbit space calU\_e/U(1) with the same quotient topology, although we do not identify its atlas with the chosen quotient-theorem atlas. Connectedness and compactness depend only on that topology. Put n=dim\_C calX-1. For the package manifold, which is an n-manifold without boundary, every local homology group H\_n(breve-calU,breve-calU\textbackslash\{\}\{x\};Z) is infinite cyclic, and its smooth orientation is a coherent choice of generators of these groups. Local homology and its continuation maps depend only on the underlying topology. Consequently, for the quotient-theorem smooth structure every point has nonzero local homology precisely in degree n; the local half-space calculation for a smooth m-manifold forces m=n and forces every point to be an interior point, so this structure also has no boundary. The same coherent local-homology generators define an orientation for it. Thus the chosen quotient-theorem structure retains all connected, compact, orientable, boundaryless, and dimension-n conclusions of the package without invoking uniqueness of the package atlas. \steptag{archived} \\[4pt]
\step{1.3.1.3} Continuity of mu and sigma implies mu(calU\_e x calU\_e) and sigma(calU\_e) are connected; they contain mu(J,J)=J and sigma(J)=J, so they lie in calU\_e. The identities mu(cU,dV)=cd mu(U,V) and sigma(cU)=conj(c)sigma(U) give p(mu(cU,dV))=p(mu(U,V)) and p(sigma(cU))=p(sigma(U)), proving that breve-mu and breve-sigma are well-defined. Since p and mu,sigma are smooth, local smooth sections of the surjective submersions p x p and p express the descended maps locally as p o mu o s\_\{12\} and p o sigma o s; representative-independence makes these expressions agree on overlaps, hence both descended maps are smooth. Finally (A\_5) gives breve-mu(breve-e,[U])=[U]=breve-mu([U],breve-e) and breve-sigma(breve-e)=breve-e. \steptag{archived} \\[4pt]
\step{1.3.2} Since epsilon\_r<=epsilon\_B\textasciicircum{}r<=epsilon\_*\textasciicircum{}r<=1/C\_grp, q\_*:=C\_grp*epsilon\_r lies in [0,1]. For U in calU\_e put U\_0:=mu(sigma(U),U), U\_1:=J. Then (A\_5) gives ||U\_0-U\_1||<=q\_*; (R) gives C\_path*q\_*<=1/4 and C\_path*(q\_*+epsilon\_r*q\_*+q\_*\textasciicircum{}2)=eta\_*<r\_-:=delta\_*-C\_pol*(epsilon\_r*delta\_*+delta\_*\textasciicircum{}2). Hence (A\_6) supplies H(t,U\_0,J). The formula F(t,[U]):=[H(t,mu(sigma(U),U),J)] is independent of the representative because sigma(cU)=conj(c)sigma(U) and mu(conj(c)sigma(U),cU)=mu(sigma(U),U). It stays in calU\_e because each path starts there; it is continuous because on every quotient neighborhood a local section s of p gives the continuous representative formula [H(t,mu(sigma(s(-)),s(-)),J)], and representative-independence makes these formulas agree. It begins at breve-mu(breve-sigma([U]),[U]) and ends at breve-e. It is basepoint-preserving because sigma(J)=J, mu(J,J)=J, and H(t,J,J)=u\_delta(J)=J by (A\_4). Thus breve-sigma is a left inversion in def-h-space-left-inversion; it is smooth by node 1.3.1. The exact unit identities make the two H-space unit homotopies constant, and breve-calU is connected by lem-stage1-quotient-manifold-package. Therefore (breve-calU,breve-mu,breve-e) is a connected H-space, and node 1.1 plus the all-calU equivariance and error estimate of node 1.2 supplies every remaining asserted conclusion. \steptag{validated} \\[6pt]
\step{1.3.2.1} Assume the declared dependencies 1.1, 1.2, and 1.3.1 have been validated. Node 1.2 supplies the all-calU identities sigma(cU)=conj(c)sigma(U), mu(cU,dV)=cd mu(U,V), and ||sigma(U)-U\textasciicircum{}dagger||<=C\_grp*epsilon\_r. Node 1.3.1 supplies the invariant component calU\_e, the quotient p:calU\_e->breve-calU with local sections, and the well-defined smooth descended maps breve-mu([U],[V])=[mu(U,V)] and breve-sigma([U])=[sigma(U)], including exact unit and basepoint identities. Since epsilon\_r<=epsilon\_B\textasciicircum{}r<=epsilon\_*\textasciicircum{}r<=1/C\_grp, q\_*:=C\_grp*epsilon\_r belongs to [0,1]. For U in calU\_e, U\_0:=mu(sigma(U),U) lies in calU\_e by node 1.3.1 and (A\_5) gives ||U\_0-J||<=q\_*; the guards in (R) give C\_path*q\_*<=1/4 and C\_path*(q\_*+epsilon\_r*q\_*+q\_*\textasciicircum{}2)=eta\_*<delta\_*-C\_pol*(epsilon\_r*delta\_*+delta\_*\textasciicircum{}2). Thus (A\_6) defines H(t,U\_0,J), joining U\_0 to J. Its image lies in calU\_e because it is a connected path in calU starting at U\_0 in the connected component calU\_e. Define F(t,[U]):=p(H(t,mu(sigma(U),U),J)). This is representative-independent: replacing U by cU and using the two equivariances from node 1.2 gives mu(sigma(cU),cU)=mu(conj(c)sigma(U),cU)=mu(sigma(U),U). It is continuous because the local sections of p from node 1.3.1 express it locally as p(H(t,mu(sigma(s(x)),s(x)),J)); joint continuity in (A\_6) and continuity of mu,sigma,p make this continuous, and representative-independence makes the local formulas agree. At t=0 it equals breve-mu(breve-sigma([U]),[U]); at t=1 it equals breve-e. At [U]=breve-e, (A\_5) gives sigma(J)=J and mu(J,J)=J, while (A\_6) gives H(t,J,J)=J (its asserted path joins J to J and its displayed formula plus (A\_4) gives u\_delta(J)=J), so F is basepoint-preserving. Hence breve-sigma is a left inversion by def-h-space-left-inversion, and it is smooth by node 1.3.1. The exact descended unit identities of node 1.3.1 give the two constant basepoint-preserving H-space unit homotopies, and node 1.1 gives connectedness and the compact orientable boundaryless manifold conclusion with dimension dim\_C(calX)-1. Therefore the connected H-space, smooth-left-inversion, manifold, equivariance, and error conclusions asserted by node 1.3.2 all follow, with no use of an undeclared or unvalidated sibling. \steptag{validated} \\[4pt]
\end{proofbox}

\end{document}
